\subsection{CHARACTERIZATION OF INTERVALS}

PROPOSITION 1. A non-empty subset A of R is an interval if and only if, whenever a and b are any two points of A such that a \textless{} b, the closed interval {[}a, b{]} is contained in A.

The condition is clearly necessary. Conversely, suppose that it is satisfied. If A is neither bounded above nor below it must be the whole of R, for if x is any point of R there are then two points a, b of A such that a \textless{} x \textless{} b. If A is bounded above but not below, let k be its least upper bound; then for any x \textless{} k there exist a and b in A such that a \textless{} x \textless{} b ≤ k, hence \(x \in A\) , and therefore A can only be one of the two intervals \(] \leftarrow, k],] \leftarrow, k[\) . The argument is similar in the other cases.

